%Data institusi

%\input{Dozentdaten}

%\renewcommand{\fachbereich}{Fachbereich}

%\renewcommand{\dozent}{Prof. Dr. . }

%Data ujian

\renewcommand{\klausurgebiet}{ }

\renewcommand{\klausurtyp}{ }

\renewcommand{\klausurdatum}{ . 20}

\klausurvorspann {\fachbereich} {\klausurdatum} {\dozent} {\klausurgebiet} {\klausurtyp}

%Data untuk tabel nilai berikut

\renewcommand{\aeins}{  3  }

\renewcommand{\azwei}{  3   }

\renewcommand{\adrei}{  6  }

\renewcommand{\avier}{  3  }

\renewcommand{\afuenf}{  3  }

\renewcommand{\asechs}{  5  }

\renewcommand{\asieben}{  13  }

\renewcommand{\aacht}{  4  }

\renewcommand{\aneun}{  3  }

\renewcommand{\azehn}{  7  }

\renewcommand{\aelf}{  6  }

\renewcommand{\azwoelf}{  3  }

\renewcommand{\adreizehn}{  6   }

\renewcommand{\avierzehn}{  65  }

\renewcommand{\afuenfzehn}{    }

\renewcommand{\asechzehn}{    }

\renewcommand{\asiebzehn}{    }

\renewcommand{\aachtzehn}{    }

\renewcommand{\aneunzehn}{    }

\renewcommand{\azwanzig}{    }

\renewcommand{\aeinundzwanzig}{    }

\renewcommand{\azweiundzwanzig}{    }

\renewcommand{\adreiundzwanzig}{    }

\renewcommand{\avierundzwanzig}{    }

\renewcommand{\afuenfundzwanzig}{    }

\renewcommand{\asechsundzwanzig}{    }

\punktetabelledreizehn

\klausurnote

\newpage

\setcounter{section}{0}

\inputaufgabegibtloesung
{3}
{

Definisikan istilah-istilah berikut
\zusatzklammer {yang dicetak miring} {} {.}
\aufzaehlungsechs{Suatu
\stichwort {ruang Hausdorff} {}
$X$.

}{Suatu
\stichwort {pemetaan diferensiabel} {}
\maabbdisp {\varphi} { L } { M
} {}
antara dua manifold diferensiabel
\mathkor {} {L} {dan} {M} {.}

}{
\stichwort {Bundel tangen} {}
dari suatu manifold diferensiabel $M$.

}{
\stichwort {Matriks fundamental kedua} {}
dari parametrisasi lokal
\maabbdisp {} {V} {U \subseteq Y
} {}
\zusatzklammer {dengan

\mavergleichskette
{\vergleichskette
{ V
}
{ \subseteq }{ \R^{n-1}
}
{  }{
}
{    }{
}
{   }{
}
}
{}{}{}
terbuka} {} {}
suatu hiperpermukaan diferensiabel

\mavergleichskette
{\vergleichskette
{ Y
}
{ \subseteq }{ \R^n
}
{  }{
}
{    }{
}
{   }{
}
}
{}{}{.}

}{
\stichwort {Turunan eksterior} {} dari suatu bentuk diferensial terdiferensialkan secara kontinu
\mathl{\omega \in
{ \mathcal E }^{ k } (  M   )}{} pada manifold diferensiabel $M$.

}{Suatu \stichwort {partisi kesatuan yang tersubordinasi pada penutup} {,} dengan

\mavergleichskette
{\vergleichskette
{  X
}
{ = }{ \bigcup_{i \in I} W_i
}
{  }{
}
{    }{
}
{   }{
}
}
{}{}{}
suatu
\definitionsverweis {penutup terbuka}{}{}
dari
\definitionsverweis {ruang topologis}{}{}
$X$.
}

}
{} {}

\inputaufgabegibtloesung
{3}
{

Nyatakan teorema-teorema berikut.
\aufzaehlungdrei{
\stichwort {Teorema tentang sifat aljabar pemetaan Weingarten} {.}}{Teorema tentang keterorientasian dan bentuk volume positif.}{
\stichwort {Teorema titik tetap Brouwer} {.}}

}
{} {}

\inputaufgabegibtloesung
{6}
{

Tentukan
\definitionsverweis {kelengkungan}{}{}
di setiap titik lingkaran satuan yang diberikan secara implisit oleh

\mavergleichskettedisp
{\vergleichskette
{ h(x,y)
}
{ =} { x^2+y^2
}
{ =} { 1
}
{ } {
}
{ } {
}
}
{}{}{}
dengan
[[Ebene differenzierbare Kurve/Faser/Krümmung/Fakt|Lemma 3.11 (Geometri Diferensial (Osnabrück 2023))]]
dan medan gradien $h$.

}
{} {}

\inputaufgabegibtloesung
{3}
{

Misalkan $X$ suatu
\definitionsverweis {ruang Hausdorff}{}{.}
Tunjukkan bahwa
\definitionsverweis {diagonal}{}{}

\mavergleichskettedisp
{\vergleichskette
{ \triangle
}
{ =} { \left\{ (x,y) \in X \times X  \mid   x = y        \right\}
}
{ } {
}
{ } {
}
{ } {
}
}
{}{}{}
merupakan
\definitionsverweis {subhimpunan tertutup}{}{}
dalam
\definitionsverweis {ruang produk}{}{}

\mathl{X \times X}{}.

}
{} {}

\inputaufgabegibtloesung
{3}
{

Tunjukkan bahwa ketiga manifold satu dimensi
\mathlistdisp {S^1 = \left\{ (x,y) \in \R^2  \mid   \betrag { (x,y) } = 1         \right\}} {} {\R} {dan} {\R \setminus \{0\}} {}
tidak ada dua di antaranya yang homeomorfik.

}
{} {}

\inputaufgabegibtloesung
{5}
{

Misalkan

\mavergleichskette
{\vergleichskette
{ G
}
{ \subseteq }{ \R^m
}
{  }{
}
{    }{
}
{   }{
}
}
{}{}{}
\definitionsverweis {terbuka}{}{}
dan misalkan

\mavergleichskette
{\vergleichskette
{ M
}
{ \subseteq }{ G
}
{  }{
}
{    }{
}
{   }{
}
}
{}{}{}
suatu
$n$-\definitionsverweis {dimensional}{}{}
\definitionsverweis {submanifold tertutup}{}{,}
yang
\definitionsverweis {berorientasi}{}{}
dan dilengkapi dengan struktur Riemann terinduksi serta
\definitionsverweis {bentuk volume kanonik}{}{}
$\omega$. Misalkan

\mavergleichskette
{\vergleichskette
{ W
}
{ \subseteq }{ \R^n
}
{  }{
}
{    }{
}
{   }{
}
}
{}{}{}
terbuka dan misalkan
\maabbdisp {\varphi} { W } { U
} {}
suatu
\definitionsverweis {difeomorfisme}{}{}
dengan himpunan terbuka

\mavergleichskette
{\vergleichskette
{ U
}
{ \subseteq }{ M
}
{  }{
}
{    }{
}
{   }{
}
}
{}{}.
Tunjukkan bahwa pada $W$ berlaku rumus

\mavergleichskettealignhandlinks
{\vergleichskettealignhandlinks
{\varphi^* \left(  \omega \vert_U  \right)
}
{ =} { \left( \det  \left( \left\langle  \begin{pmatrix}  { \frac{   \partial  \varphi_1   }{  \partial   x_i   } }  \\ \vdots \\  { \frac{   \partial  \varphi_m   }{  \partial   x_i   } }   \end{pmatrix}  ,  \begin{pmatrix}  { \frac{   \partial  \varphi_1   }{  \partial   x_j   } }  \\ \vdots \\  { \frac{   \partial  \varphi_m   }{  \partial   x_j   } }   \end{pmatrix}   \right\rangle \right)_{1 \leq i,j \leq n} \right)^{1/2}  dx_1 \wedge \ldots \wedge dx_n
}
{ =} { \left( \det  \left(  { \frac{   \partial  \varphi_1   }{  \partial   x_i   } } { \frac{   \partial  \varphi_1   }{  \partial   x_j   } }  + \cdots +  { \frac{   \partial  \varphi_m   }{  \partial   x_i   } } { \frac{   \partial  \varphi_m   }{  \partial   x_j   } }  \right)_{1 \leq i,j \leq n} \right)^{1/2}  dx_1 \wedge \ldots \wedge dx_n
}
{ } {
}
{ } {
}
}
{}
{}{}{.}

}
{} {}

\inputaufgabegibtloesung
{13 (2+3+2+2+2+2)}
{

Kita tinjau

\mavergleichskette
{\vergleichskette
{ \R^6
}
{ = }{ \left(  \R^2  \right)^3
}
{  }{
}
{    }{
}
{   }{
}
}
{}{}{}
sebagai himpunan semua segitiga
\zusatzklammer {termasuk yang degenerat} {} {}
dengan mengidentifikasi segitiga yang mempunyai titik-titik sudut
\zusatzklammer {terurut} {} {}

\mavergleichskette
{\vergleichskette
{ P_1
}
{ = }{ \begin{pmatrix}  x_1  \\ y_1    \end{pmatrix}
}
{  }{
}
{    }{
}
{   }{
}
}
{}{}{,}

\mavergleichskette
{\vergleichskette
{ P_2
}
{ = }{ \begin{pmatrix}  x_2  \\ y_2    \end{pmatrix}
}
{  }{
}
{    }{
}
{   }{
}
}
{}{}{}
dan

\mavergleichskette
{\vergleichskette
{ P_3
}
{ = }{ \begin{pmatrix}  x_3  \\ y_3    \end{pmatrix}
}
{  }{
}
{    }{
}
{   }{
}
}
{}{}{,}
dengan tupel koordinat

\mathdisp {(x_1,y_1,x_2,y_2,x_3,y_3)} {  }

.
\aufzaehlungsechs{Tunjukkan bahwa himpunan segitiga dengan dua titik sudut berimpit merupakan subhimpunan tertutup $\R^6$
\zusatzklammer {jadi komplemennya merupakan himpunan terbuka dalam $\R^6$ yang kita namai $U$} {} {.}
}{Tunjukkan bahwa himpunan segitiga dengan ketiga titik sudut segaris merupakan subhimpunan tertutup $\R^6$
\zusatzklammer {jadi komplemennya, yang terdiri atas semua segitiga tak degenerat, merupakan himpunan terbuka dalam $\R^6$ yang kita namai $V$} {} {.}
}{Buat aturan fungsi yang mendeskripsikan pemetaan
\maabbdisp {F} {\R^6} {\R
} {}
yang memetakan suatu segitiga ke kelilingnya.
}{Tunjukkan bahwa fungsi $F$ dari bagian (3) terdiferensialkan secara kontinu pada himpunan $U$.
}{Hitung turunan parsial $F$ terhadap $x_1$ pada $U$.
}{Tunjukkan bahwa himpunan segitiga tak degenerat berkeliling $1$ membentuk
\definitionsverweis {submanifold tertutup}{}{}
dari $V$. Berapakah dimensinya?
}

}
{} {}

\inputaufgabegibtloesung
{4}
{

Misalkan
\mathkor {} {r} {dan} {R} {}
bilangan-bilangan real positif dengan

\mavergleichskette
{\vergleichskette
{ 0
}
{ < }{ r
}
{ < }{ R
}
{    }{
}
{   }{
}
}
{}{}{,}
dan tinjau torus
\zusatzklammer {tertanam} {} {}

\mavergleichskettedisp
{\vergleichskette
{ T
}
{ =} { \left\{ (x,y,z) \in \R^3  \mid   \left(  \sqrt{x^2+y^2} -R  \right)^2 +z^2  = r^2         \right\}
}
{ \subseteq} { \R^3
}
{ } {
}
{ } {
}
}
{}{}{}
dengan
\definitionsverweis {struktur Riemann terinduksi}{}{.}
Tunjukkan bahwa untuk setiap

\mavergleichskette
{\vergleichskette
{ \alpha
}
{ \in }{ [0,2 \pi [
}
{  }{
}
{    }{
}
{   }{
}
}
{}{}{}
pemetaan
\maabbeledisp {\Psi_\alpha} {\R^3} {\R^3
} { \left(  x   , \,   y  , \,  z                 \right) } { \left(  x \cos \alpha - y \sin  \alpha   , \,   x \sin \alpha + y \cos  \alpha  , \,  z                 \right)
} {,}
menginduksi suatu
\definitionsverweis {isometri}{}{}
dari $T$ ke dirinya sendiri.

}
{} {}

\inputaufgabe
{3}
{

Deskripsikan suatu model untuk bidang hiperbolik.

}
{} {}

\inputaufgabegibtloesung
{7}
{

Buktikan Theorema egregium.

}
{} {}

\inputaufgabegibtloesung
{6 (1+2+2+1)}
{

Kita tinjau
\definitionsverweis {bentuk diferensial}{}{}

\mavergleichskettedisp
{\vergleichskette
{ \omega
}
{ =} { ydx+zdy +xdz
}
{ } {
}
{ } {
}
{ } {
}
}
{}{}{}
pada $\R^3$ dan pemetaan
\maabbeledisp {\varphi} {\R^2} {\R^3
} { \left(  u   , \,  v                   \right) } { \left( u^2  , \,  v^2  , \,  uv                 \right)
} {.}
\aufzaehlungvier{Hitung turunan eksterior $\omega$.
}{Hitung tarikan balik $\varphi^* \omega$ dari $\omega$ di bawah $\varphi$.
}{Hitung turunan eksterior $\varphi^* \omega$ pada $\R^2$.
}{Hitung tarikan balik $\varphi^* d \omega$ dari $d\omega$ di bawah $\varphi$ secara independen dari bagian (3).
}

}
{} {}

\inputaufgabegibtloesung
{3}
{

Misalkan $M$ suatu
\definitionsverweis {manifold diferensiabel}{}{}
dan $\omega$ suatu
\definitionsverweis {bentuk diferensial}{}{}
yang
\definitionsverweis {eksak}{}{}
pada $M$. Misalkan $S$ suatu
\definitionsverweis {manifold diferensiabel}{}{}
yang
\definitionsverweis {kompak}{}{}
dan
\definitionsverweis {berorientasi}{}{}
\zusatzklammer {tanpa batas} {} {}
dengan
\definitionsverweis {basis terhitung bagi topologinya}{}{}
dan misalkan
\maabbdisp {\varphi} { S } { M
} {}
suatu
\definitionsverweis {pemetaan terdiferensialkan secara kontinu}{}{.}
Tunjukkan bahwa

\mavergleichskettedisp
{\vergleichskette
{
\int_{  S  }   \varphi^*\omega
}
{ =} { 0
}
{ } {
}
{ } {
}
{ } {
}
}
{}{}{.}

}
{} {}

\inputaufgabegibtloesung
{6}
{

Buktikan teorema tentang retraksi ke batas pada manifold.

}
{} {}

 [[Kategorie:Latexseite]]
